%% sections/datasheet.tex  (APPENDIX)
%% \owner{DOCS subagent}  GOAL: a "Datasheet for Datasets" (Gebru et al.) for the dev split.
%% Most answers already exist in bench/dataset.py, bench/data/manifest.json, and bench/README.md
%% — assemble, do not invent. Keep each answer short. Sections (Gebru standard):
%%   Motivation; Composition; Collection/Generation Process; Preprocessing/Labelling; Uses;
%;   Distribution; Maintenance. ALSO (DOCS subagent's other deliverables, separate files —
%;   NOT this .tex): refresh root CITATION.cff + codemeta.json, add a machine-readable
%;   Croissant record (croissant.json) for the dev split, and a top-level DATASHEET.md
%;   mirroring this appendix. Confirm LICENSE (MIT) is present.

\section{Datasheet}
\label{app:datasheet}

We follow the ``Datasheets for Datasets'' template of \citet{gebru2021datasheets}.
The answers describe the public \texttt{dev} split frozen in \texttt{bench/data/dev/}
and indexed by \texttt{bench/data/manifest.json} (dataset version \texttt{2.0}).

\paragraph{Motivation.}
The dataset was created to evaluate and train code models under conditions that
mined corpora cannot provide: a \emph{known-optimal reference} to grade against,
two \emph{orthogonal} difficulty axes that can be moved independently, and a
\emph{contamination control}. It supports three tasks on a single shared label set
--- refactoring (rewrite the spaghetti), LLM-as-judge code-quality ranking, and
output-prediction comprehension. It was created by the author of \tool{} (see
\texttt{CITATION.cff}); no external funding is associated with this release.

\paragraph{Composition.}
Each instance is one program rendered from a clean JSON intermediate representation
(IR) built from four operation primitives (\texttt{KEY\_VALUE\_LOOKUP},
\texttt{MEMBERSHIP\_CHECK}, \texttt{AGGREGATE}, \texttt{CONDITIONAL\_SELECT}). An
instance is located on four axes: \emph{incidental} complexity --- a strictly-nested
6-level messiness knob at identical semantics, with by-construction maintainability
rank \texttt{clean} $\subset$ \texttt{minimal} $\subset$ \texttt{light} $\subset$
\texttt{standard} $\subset$ \texttt{heavy} $\subset$ \texttt{max} (\texttt{clean} is
the known-optimal baseline); \emph{intrinsic} complexity --- per-family scale knobs
(\texttt{config\_resolver}~$N$, \texttt{allowlist}~$L$, \texttt{agg\_stats}~$W$,
\texttt{threshold\_select}~$T$) and operation count; \emph{language} ---
\{python, javascript, go, java, cpp\}; and input \emph{variant} ---
\texttt{v0..v4} perturbations used for differential grading. The public \texttt{dev}
split has 100 items (50 base + 50 variant) across seven families:
\texttt{config\_resolver}, \texttt{allowlist}, \texttt{status\_router},
\texttt{discovery\_pipeline}, \texttt{fsm\_transition}, \texttt{agg\_stats},
\texttt{threshold\_select}. Each item ships its IR, the rendered sources for all five
languages at every profile, and ground truth (oracle output, the known-optimal clean
baseline, knob coordinates, and the planner's enabled \texttt{SPAGH\_*} set). A second,
\textbf{private} held-out \texttt{test} split exists but is not distributed: it is
minted on demand from \texttt{BENCH\_HELDOUT\_SEED} and is never committed. The release
embeds a unique canary GUID
(\texttt{spaghetti-architect-bench:\allowbreak canary:\allowbreak 9f1d4e7a-3b6c-\allowbreak 4f20-bd8e-\allowbreak 1a2c3d4e5f60}) in
\texttt{bench/data/CANARY.txt} and the manifest. The data are wholly synthetic: no
personal, sensitive, or human-subject data, and no offensive content.

\paragraph{Collection / Generation process.}
There is no collection: every instance is generated. A single IR generator
(\texttt{eval.gen\_samples.build(extended=True)}; there is deliberately no second
sampler) draws all RNG-sourced literals from the public seed \texttt{20260619}; the
project transpiler (\texttt{src/} \texttt{Engine}/\texttt{Planner}) renders the five
language targets at each profile; and an oracle validator computes the reference
output. Every minted IR is checked with the real parser
(\texttt{src.nodes.parser.parse}) before it enters the split, so a structurally
invalid sample can never be included. Generation is byte-deterministic given the
seed. The private \texttt{test} split re-mints numeric literals from a private seed
and applies a structure-preserving, one-way (SHA-256) string salt (novelty Tier~A),
adding held-out structural shapes (Tier~B) and out-of-distribution compositions
(Tier~C); the held-out seed is scrubbed to \texttt{PRIVATE} and never written to any
artifact.

\paragraph{Preprocessing / Labelling.}
Labels are correct \emph{by construction}, not human-annotated: the oracle output is
computed from the IR; the incidental rank is the strictly-nested \texttt{SPAGH\_*}
inclusion chain; the intrinsic coordinates are the generator's own knob settings.
No cleaning, filtering, or external labelling is applied, and no instances are
discarded after validation. The raw IR is retained alongside the rendered sources in
every \texttt{dev} record, so the full provenance of each instance is reproducible.

\paragraph{Uses.}
The dataset is intended for machine-learning research on code: (i)
\emph{contamination-controlled LLM evaluation} --- score on \texttt{dev} for
reproducibility, then on the private \texttt{test} split, graded by novelty tier, to
measure the memorization gap; (ii) \emph{training and probing} of code models on a
controlled, by-construction-labelled signal (refactoring targets, quality ranking,
output prediction); and (iii) \emph{robustness studies} that vary incidental vs
intrinsic complexity independently to attribute failures to presentation mess rather
than problem size. The shipped labels should be treated as a controlled stimulus,
not as a model of natural-corpus difficulty; the \texttt{dev} split is public and so,
once released, must be assumed potentially in future training sets --- only the
private \texttt{test} split provides a contamination guarantee. It should not be used
to claim that any model is ``contamination-proof''; the defense is the minting
protocol, not the artifact.

\paragraph{Distribution.}
The dataset is distributed with the tool on GitHub
(\url{https://github.com/KurathSec/Spaghetti-Architect}) under the MIT license, and
is also described by a machine-readable MLCommons Croissant record
(\texttt{croissant.json}) and a top-level \texttt{DATASHEET.md}. The public
\texttt{dev} split is included in the repository; the \texttt{test} split is private
and withheld by design. A persistent archival DOI (Zenodo) will be assigned to the
tagged release.

\paragraph{Maintenance.}
The dataset is maintained by the author via the GitHub repository (issues / pull
requests); contact is the corresponding author in \texttt{CITATION.cff}. It is
versioned (\texttt{dataset\_version} in the manifest; current \texttt{2.0}) and is
fully regenerable with \texttt{python3 bench/dataset.py --freeze}. The generation
contract is the extension point: new families are added by extending the single IR
generator and re-freezing, which keeps labels by-construction and the splits
reproducible. The private held-out seed can be rotated to refresh the \texttt{test}
split without changing the public release.

\paragraph{Author statement and licensing.}
The author confirms that the generator and the public \texttt{dev} split are released
under the MIT License (see \texttt{LICENSE}), that the data are wholly synthetic and
infringe no third-party rights, and that the author bears all responsibility in case of
any violation of rights. Hosting, licensing, and long-term access follow the plan above:
the artifact is served from the public GitHub repository and will be archived under a
persistent Zenodo DOI (to be assigned at release); the project is maintained by the author
through the repository's issue and pull-request channels.
